% GENERATED FILE. Edit canonical lessons, then run npm run generate:books.
% canonical-source-hash: fnv1a64:f202c9b261b9ae86
% canonical-lessons: KA-C345-adagu, KA-C345-talemaresiko, KA-C345-bennattu, KA-C345-gudhacara, KA-C345-nusulu

\chapter{The Spy in the Palace}
\label{ch:ka-the-spy-in-the-palace}
\hlchaptermodality{345}

\begin{chapteropening}
\textbf{By the end of this chapter:} \emph{I can say to hide oneself, to go into hiding, to chase, a spy and to sneak in.}

Complete the last lesson of chapter 345: I can say to hide oneself, to go into hiding, to chase, a spy and to sneak in.
\end{chapteropening}

\section[\texorpdfstring{\kn{ಅಡಗು} (aḍagu)}{ಅಡಗು (aḍagu)}]{\kn{ಅಡಗು} (aḍagu) — to hide oneself}

\label{lesson:KA-C345-adagu}



\begin{quote}
Before the new one: say the Kannada for a chariot, then the Kannada for to do battle.
\end{quote}

\subsection*{You'll want to know: \kn{ಅಡಗು}}
\textbf{\kn{ಅಡಗು}} — \emph{aḍagu} — \textquotedblleft{}to hide oneself\textquotedblright{}.

\begin{quote}
\kn{ಕಳ್ಳ} \kn{ಬಾವಿಯ} \kn{ಹಿಂದೆ} \kn{ಅಡಗಿದನು}.
\end{quote}

\begin{quote}
\emph{kaḷḷa bāviya hinde aḍagidanu.}
\end{quote}

\begin{quote}
The thief hid behind the well.
\end{quote}

Past (I …): \kn{ಅಡಗಿದೆ}, \emph{aḍagide}.

To hide oneself is \kn{ಅಡಗು}; to hide something is \kn{ಅಡಗಿಸು} or \kn{ಬಚ್ಚಿಡು}.

\begin{cousinweb}
\textbf{\kn{ಅಡಗು}} (\emph{aḍagu}) is native Dravidian: DEDR 63 lists Kannada \emph{aḍagu} beside Tamil \emph{aṭaṅku}, \textquotedblleft{}to be compressed, to cease, to disappear\textquotedblright{}.
\end{cousinweb}

\begin{grammarlens}[title={What you've built}]
One more word for telling what happened.
\end{grammarlens}

\subsection*{Your turn}
\begin{itemize}
\raggedright
  \item \emph{Say it:} \emph{aḍagu}
  \item \emph{Say it:} \emph{aḍagu}, once more
  \item \emph{Say it:} \emph{kaḷḷa bāviya hinde aḍagidanu.} — The thief hid behind the well.
  \item \emph{Recall:} say the Kannada for a chariot, then the Kannada for to do battle, then say \emph{aḍagu} again
\end{itemize}

\begin{tcolorbox}[breakable,colback=teal!4,colframe=teal!35!black,title={Before you move on}]
What does \kn{ಅಡಗು} mean? (\textquotedblleft{}To hide oneself\textquotedblright{}.) Say it once more.
\end{tcolorbox}
\section[\texorpdfstring{\kn{ತಲೆಮರೆಸಿಕೊ} (talemaresiko)}{ತಲೆಮರೆಸಿಕೊ (talemaresiko)}]{\kn{ತಲೆಮರೆಸಿಕೊ} (talemaresiko) — to go into hiding}

\label{lesson:KA-C345-talemaresiko}



\begin{quote}
Before the new one: say the Kannada for to do battle, then the Kannada for to hide oneself.
\end{quote}

\subsection*{You'll want to know: \kn{ತಲೆಮರೆಸಿಕೊ}}
\textbf{\kn{ತಲೆಮರೆಸಿಕೊ}} — \emph{talemaresiko} — \textquotedblleft{}to go into hiding\textquotedblright{}.

\begin{quote}
\kn{ಆರೋಪಿ} \kn{ಆರು} \kn{ತಿಂಗಳು} \kn{ತಲೆಮರೆಸಿಕೊಂಡಿದ್ದನು}.
\end{quote}

\begin{quote}
\emph{ārōpi āru tiṅgaḷu talemaresikoṇḍiddanu.}
\end{quote}

\begin{quote}
The accused had been in hiding for six months.
\end{quote}

Past (I …): \kn{ತಲೆಮರೆಸಿಕೊಂಡೆ}, \emph{talemaresikoṇḍe}.

Literally \textquotedblleft{}to hide one's head\textquotedblright{}: a news phrase for someone wanted by the police.

\begin{cousinweb}
\textbf{\kn{ತಲೆಮರೆಸಿಕೊ}} (\emph{talemaresiko}) is a native Dravidian compound: \emph{tale}, \textquotedblleft{}head\textquotedblright{} (DEDR 3103), and \emph{maresu}, \textquotedblleft{}to hide\textquotedblright{}, from \emph{mare} (DEDR 4760), with -\emph{ko}, from \emph{koḷ}, \textquotedblleft{}to take\textquotedblright{} (DEDR 2151).
\end{cousinweb}

\begin{grammarlens}[title={What you've built}]
One more word for telling what happened.
\end{grammarlens}

\subsection*{Your turn}
\begin{itemize}
\raggedright
  \item \emph{Say it:} \emph{talemaresiko}
  \item \emph{Say it:} \emph{talemaresiko}, once more
  \item \emph{Say it:} \emph{ārōpi āru tiṅgaḷu talemaresikoṇḍiddanu.} — The accused had been in hiding for six months.
  \item \emph{Recall:} say the Kannada for to do battle, then the Kannada for to hide oneself, then say \emph{talemaresiko} again
\end{itemize}

\begin{tcolorbox}[breakable,colback=teal!4,colframe=teal!35!black,title={Before you move on}]
What does \kn{ತಲೆಮರೆಸಿಕೊ} mean? (\textquotedblleft{}To go into hiding\textquotedblright{}.) Say it once more.
\end{tcolorbox}
\section[\texorpdfstring{\kn{ಬೆನ್ನಟ್ಟು} (bennaṭṭu)}{ಬೆನ್ನಟ್ಟು (bennaṭṭu)}]{\kn{ಬೆನ್ನಟ್ಟು} (bennaṭṭu) — to chase}

\label{lesson:KA-C345-bennattu}



\begin{quote}
Before the new one: say the Kannada for to hide oneself, then the Kannada for to go into hiding.
\end{quote}

\subsection*{You'll want to know: \kn{ಬೆನ್ನಟ್ಟು}}
\textbf{\kn{ಬೆನ್ನಟ್ಟು}} — \emph{bennaṭṭu} — \textquotedblleft{}to chase\textquotedblright{}.

\begin{quote}
\kn{ಪೊಲೀಸರು} \kn{ಕಳ್ಳನ} \kn{ಕಾರನ್ನು} \kn{ಬೆನ್ನಟ್ಟಿದರು}.
\end{quote}

\begin{quote}
\emph{polīsaru kaḷḷana kārannu bennaṭṭidaru.}
\end{quote}

\begin{quote}
The police chased the thief's car.
\end{quote}

Past (I …): \kn{ಬೆನ್ನಟ್ಟಿದೆ}, \emph{bennaṭṭide}.

\kn{ಬೆನ್ನು}, back, and \kn{ಅಟ್ಟು}, to run after: to run at someone's back.

\begin{cousinweb}
\textbf{\kn{ಬೆನ್ನಟ್ಟು}} (\emph{bennaṭṭu}) is a native Dravidian compound: \emph{bennu}, \textquotedblleft{}the back\textquotedblright{} (DEDR 5488), and \emph{aṭṭu}, \textquotedblleft{}to follow, to run after, to pursue\textquotedblright{}, which DEDR 79 lists for Kannada.
\end{cousinweb}

\begin{grammarlens}[title={What you've built}]
One more word for telling what happened.
\end{grammarlens}

\subsection*{Your turn}
\begin{itemize}
\raggedright
  \item \emph{Say it:} \emph{bennaṭṭu}
  \item \emph{Say it:} \emph{bennaṭṭu}, once more
  \item \emph{Say it:} \emph{polīsaru kaḷḷana kārannu bennaṭṭidaru.} — The police chased the thief's car.
  \item \emph{Recall:} say the Kannada for to hide oneself, then the Kannada for to go into hiding, then say \emph{bennaṭṭu} again
\end{itemize}

\begin{tcolorbox}[breakable,colback=teal!4,colframe=teal!35!black,title={Before you move on}]
What does \kn{ಬೆನ್ನಟ್ಟು} mean? (\textquotedblleft{}To chase\textquotedblright{}.) Say it once more.
\end{tcolorbox}
\section[\texorpdfstring{\kn{ಗೂಢಚಾರ} (gūḍhacāra)}{ಗೂಢಚಾರ (gūḍhacāra)}]{\kn{ಗೂಢಚಾರ} (gūḍhacāra) — a spy}

\label{lesson:KA-C345-gudhacara}



\begin{quote}
Before the new one: say the Kannada for to go into hiding, then the Kannada for to chase.
\end{quote}

\subsection*{You'll want to know: \kn{ಗೂಢಚಾರ}}
\textbf{\kn{ಗೂಢಚಾರ}} — \emph{gūḍhacāra} — \textquotedblleft{}a spy\textquotedblright{}.

\begin{quote}
\kn{ರಾಜನ} \kn{ಗೂಢಚಾರ} \kn{ಪಕ್ಕದ} \kn{ರಾಜ್ಯಕ್ಕೆ} \kn{ಹೋದನು}.
\end{quote}

\begin{quote}
\emph{rājana gūḍhacāra pakkada rājyakke hōdanu.}
\end{quote}

\begin{quote}
The king's spy went to the neighbouring kingdom.
\end{quote}

Plural: \kn{ಗೂಢಚಾರರು}, \emph{gūḍhacāraru}.

\kn{ಗೂಢ}, hidden, and \kn{ಚಾರ}, a goer.

\begin{cousinweb}
\textbf{\kn{ಗೂಢಚಾರ}} (\emph{gūḍhacāra}) is borrowed from Sanskrit \emph{gūḍhacāra}, \textquotedblleft{}a spy\textquotedblright{} (Monier-Williams), from \emph{gūḍha}, \textquotedblleft{}hidden\textquotedblright{}, and \emph{cāra}, \textquotedblleft{}going, a goer\textquotedblright{}.
\end{cousinweb}

\begin{grammarlens}[title={What you've built}]
One more word for telling what happened.
\end{grammarlens}

\subsection*{Your turn}
\begin{itemize}
\raggedright
  \item \emph{Say it:} \emph{gūḍhacāra}
  \item \emph{Say it:} \emph{gūḍhacāra}, once more
  \item \emph{Say it:} \emph{rājana gūḍhacāra pakkada rājyakke hōdanu.} — The king's spy went to the neighbouring kingdom.
  \item \emph{Recall:} say the Kannada for to go into hiding, then the Kannada for to chase, then say \emph{gūḍhacāra} again
\end{itemize}

\begin{tcolorbox}[breakable,colback=teal!4,colframe=teal!35!black,title={Before you move on}]
What does \kn{ಗೂಢಚಾರ} mean? (\textquotedblleft{}A spy\textquotedblright{}.) Say it once more.
\end{tcolorbox}
\section[\texorpdfstring{\kn{ನುಸುಳು} (nusuḷu)}{ನುಸುಳು (nusuḷu)}]{\kn{ನುಸುಳು} (nusuḷu) — to sneak in; to slip through}

\label{lesson:KA-C345-nusulu}



\begin{quote}
Before the new one: say the Kannada for to chase, then the Kannada for a spy.
\end{quote}

\subsection*{You'll want to know: \kn{ನುಸುಳು}}
\textbf{\kn{ನುಸುಳು}} — \emph{nusuḷu} — \textquotedblleft{}to sneak in; to slip through\textquotedblright{}.

\begin{quote}
\kn{ಗೂಢಚಾರ} \kn{ರಾತ್ರಿ} \kn{ಅರಮನೆಯೊಳಗೆ} \kn{ನುಸುಳಿದನು}.
\end{quote}

\begin{quote}
\emph{gūḍhacāra rātri aramaneyoḷage nusuḷidanu.}
\end{quote}

\begin{quote}
The spy slipped into the palace at night.
\end{quote}

Past (I …): \kn{ನುಸುಳಿದೆ}, \emph{nusuḷide}.

Also used of crossing a border secretly.

\begin{cousinweb}
\textbf{\kn{ನುಸುಳು}} (\emph{nusuḷu}) is native Dravidian: DEDR 3714 lists Kannada \emph{nusuḷ} beside Tamil words meaning \textquotedblleft{}to creep through a narrow passage, to penetrate\textquotedblright{}.
\end{cousinweb}

\begin{grammarlens}[title={What you've built}]
One more word for telling what happened.
\end{grammarlens}

\subsection*{Your turn}
\begin{itemize}
\raggedright
  \item \emph{Say it:} \emph{nusuḷu}
  \item \emph{Say it:} \emph{nusuḷu}, once more
  \item \emph{Say it:} \emph{gūḍhacāra rātri aramaneyoḷage nusuḷidanu.} — The spy slipped into the palace at night.
  \item \emph{Recall:} say the Kannada for to chase, then the Kannada for a spy, then say \emph{nusuḷu} again
\end{itemize}

\begin{tcolorbox}[breakable,colback=teal!4,colframe=teal!35!black,title={Before you move on}]
What does \kn{ನುಸುಳು} mean? (\textquotedblleft{}To sneak in; to slip through\textquotedblright{}.) Say it once more.
\end{tcolorbox}
